

\section{Introduction}
Large Language Models (\llm{s}) have rapidly progressed from simple code autocompletion~\cite{austin2021program,chen2021codex,li2022competition, xia2022less} to interactive agents capable of navigating repositories, running tests, and submitting patches end‑to‑end~\cite{yang2024sweagent,wang2024openhands,zhang2024autocoderover,xia2024agentless,carbonneaux2025cwm, liu2024large}.
Early conversational repair systems leveraged environment feedback to iteratively refine candidate fixes~\cite{xia2024automated}, while subsequent agentic frameworks (such as \sweagent~\cite{yang2024sweagent} and \openhands~\cite{wang2024openhands}) augmented LLMs with terminals, editors, and search, enabling multi‑step tool use on complex repositories. At the same time, complementary ``more agentless'' pipelines (e.g., Moatless~\cite{moatless} and \agentless~\cite{xia2024agentless}) argue that much of the perceived complexity in agent scaffold design can be replaced by dedicated prompting and workflow design.





Despite this progress, most existing agents have fixed designs and are limited in their static action space, where the scaffold implementation (including tools) is preset even when a task could benefit from more customization. In addition, manually designing an optimal software agent scaffold is extremely challenging and costly due to the infinite design space.
As a result, more recently, the community has begun to explore self-evolving software agents~\cite{zhang2025darwin,robeyns2025sica,wang2025huxley}, which iteratively modify their own scaffold implementations and empirically validate each change using offline evaluation signals from coding benchmarks.
However, these approaches add significant additional cost.
For example, a single run of \dgm on \swebench{} {costs} around \$22,000 according to the original paper~\cite{zhang2025darwin}.
Furthermore, they rely heavily on offline evolution, where improvements are typically learned on certain benchmarks and then baked into a static agent.
In this way, the learned agents may become specialized to the given benchmarks and underlying LLMs, and may not generalize well beyond that.


To bridge this gap,
we introduce \livesweagent, the first
\emph{live}, runtime self-evolving software engineering agent that expands and
revises its own capabilities \emph{on the fly} while working on a real-world issue.
Our key insight is that software agents are themselves software systems, and modern LLM-based software agents already possess the intrinsic capability to extend or modify their own implementation at runtime.
{While our idea of enabling on-the-fly self-evolution applies to all parts of the agent scaffold implementation, as a first step, we focus primarily on tool creation, as it is one of the most essential parts of a software agent.}
{\tech} starts from a simplistic agent with only bash tool access (e.g., mini-SWE-agent~\cite{yang2024sweagent}).
During {the} regular issue-solving loop, {the agent} can synthesize, modify, and execute custom tools such as editors, code search utilities, and domain-specific analyzers.
A lightweight step‑reflection prompt repeatedly asks the agent whether creating or revising a tool would accelerate progress, thereby turning tooling into a first‑class decision alongside ordinary actions like running tests.
This mechanism leaves the underlying scaffold unchanged, requires no offline training, and is agnostic to the underlying LLM.
By elevating tool creation to an explicit, iterative decision point, we unlock this hidden \emph{on-the-fly} self-improvement ability.

\livesweagent addresses the limitations of existing research.
First, with tool creation, the agent action space adapts to the problem at hand, producing task-relevant tools that precisely capture the needs to accomplish the current task.
Furthermore, by shifting improvement from offline training to online evolution, \livesweagent alleviates tedious scaffold engineering because new abilities emerge from the encountered issues themselves.
Importantly, tool synthesis is not a one-shot preprocessing action but a continuous iterative process interleaved with problem solving. The agent can refine tools as their understanding of the failure mode evolves.
Despite its minimal and simplistic design, \livesweagent is generalizable to different agent scaffolds and LLMs, with state-of-the-art open results on software issue solving.

We have evaluated \livesweagent on the widely used \sbv benchmark~\cite{sbv} and the more challenging \swebenchpro~\cite{deng2025swe}. Without any test‑time scaling, \livesweagent achieves \textbf{77.4\%} resolve rate on Verified and \textbf{45.8\%} on Pro, surpassing state-of-the-art open‑source baselines and{ even} the best commercial agent systems.
Our ablations and tooling analyses reveal that (1) custom tool creation materially improves effectiveness (higher solve rates) with minimal overhead,
(2) benefits persist across different state-of-the-art \llm backends, with better results on stronger models, demonstrating its promising future as LLMs' capabilities rapidly evolve,
and (3) synthesized tools include both generic utilities and task-aligned specializations that benefit issue-specific problem-solving.

In summary, we make the following contributions:

\begin{itemize}[leftmargin=*]
\item \textbf{The first live software agent.}  We present \livesweagent, the first \emph{live} software agent that can autonomously self-evolve its own scaffold implementation on the fly when solving real-world issues,
without any offline training or extra pipelines.
\item \textbf{Minimal and general implementation.} Our current implementation adopts a minimal and general design, where the agent starts from a minimal bash-only scaffold (e.g., \minisweagent) and self-improves on the fly by creating general or customized tools. The design is compatible with any existing software agent loop or LLM with negligible overhead, has been publicly available at: {\color{purple}\url{https://github.com/OpenAutoCoder/live-swe-agent}}.
\item \textbf{State-of-the-art performance.}
On \sbv and \swebenchpro, \livesweagent reaches 77.4\% and 45.8\% solve rates (without any test-time scaling), respectively, outperforming all existing open-source agents and commercial systems at the time of writing. To our knowledge, our \swebenchpro result is also the best reported to date.
\item \textbf{Comprehensive analysis.} We include detailed studies of when and why on-the-fly tool creation helps, how it improves efficiency, and how it compares to other designs like fixed-tool agents, workflow-based systems, and offline self-improving methods. Notably, compared to existing self-improving agents, we achieve much better performance (65.0\% solve rate on the \sbv-60 subset vs. \dgm's 53.3\%)
while being significantly less costly (no offline cost).

\item \textbf{Unified leaderboard.} For software tasks, recent LLMs are often benchmarked using manually engineered, proprietary agent scaffolds, making fair model comparison hard. 
 \livesweagent offers an open, unified, and powerful scaffold that enables genuinely fair, apples-to-apples comparison for future model releases. We are maintaining a leaderboard of recent models evaluated on real-world software tasks using \livesweagent at: {\color{purple}\url{http://live-swe-agent.github.io}}
\end{itemize}

